El trabajo ha demostrado que los objetivos planteados son alcanzables con las tecnologías actuales de desarrollo web y sin necesidad de infraestructura de servidor propia. Los resultados obtenidos permiten extraer las siguientes conclusiones.

En primer lugar, la arquitectura Next.js App Router con React Server Components ha demostrado ser particularmente adecuada para aplicaciones de este tipo, que combinan una gran cantidad de datos externos con la necesidad de interactividad en cliente. La separación entre la obtención de datos (servidor) y la presentación interactiva (cliente) simplifica el flujo de datos y reduce las fuentes de error. Los Server Components permiten realizar \textit{fetching} paralelo de múltiples APIs en el servidor sin coste para el cliente, lo que resulta en tiempos de carga significativamente inferiores a los de una SPA equivalente.

En segundo lugar, la integración de OAuth 2.0 mediante API Routes de Next.js ha probado ser una estrategia robusta y segura sin necesidad de librerías de autenticación genéricas. El patrón de almacenar tokens en cookies \texttt{httpOnly} y delegar toda la lógica de autenticación al servidor cumple con las recomendaciones de seguridad del RFC 6749 \parencite{Hardt:2012} y las guías de OWASP \parencite{OWASP:2024} para aplicaciones OAuth.

En tercer lugar, la estrategia de caché diferenciada por tipo de dato ha demostrado ser eficaz en todas sus capas: ISR en servidor para el catálogo, \texttt{no-store} para los datos de usuario protegidos por autenticación, y una caché en cliente de tres niveles para las puntuaciones externas de IMDb, Trakt y OMDb en las vistas de listas, reduciendo drásticamente el número de llamadas a APIs externas.

Finalmente, desde el punto de vista del diseño, el uso de Framer Motion para las animaciones y el sistema de superficies semi-transparentes con desenfoque de fondo ha producido una interfaz visualmente diferenciada respecto a las plataformas de \textit{tracking} existentes. Las métricas de rendimiento obtenidas son completamente compatibles con un diseño visualmente rico, siempre que las animaciones se implementen sobre la GPU (propiedades \texttt{transform} y \texttt{opacity}) y no fuercen la recarga del DOM.

\section{Discusión}

El trabajo presenta algunas limitaciones que deben reconocerse.

\textbf{Ausencia de tests automatizados.} El proyecto no incluyó una suite de tests unitarios ni de integración durante el desarrollo. La validación se realizó exclusivamente mediante pruebas manuales. Los componentes con lógica de negocio compleja (cálculo de progreso de episodios, normalización de fechas para Trakt, estrategia de selección de \textit{backdrops}) serían candidatos prioritarios para la cobertura de tests unitarios.

\textbf{Dependencia de APIs de terceros.} La plataforma depende de la continuidad y estabilidad de cinco APIs externas (TMDb, Trakt.tv, OMDb, JustWatch y Plex), ninguna de las cuales está bajo el control del proyecto. Esta es una limitación estructural de la arquitectura elegida, aceptada conscientemente por las ventajas que ofrece frente a la alternativa de construir una base de datos propia.

\textbf{Accesibilidad parcial.} La puntuación de accesibilidad en Lighthouse (88/100) indica que, aunque se han aplicado prácticas básicas, no se ha alcanzado el nivel WCAG 2.1 AA completo \parencite{W3C:2023}. En particular, la navegación exclusiva por teclado en los modales de episodios requiere trabajo adicional.

\section{Líneas de investigación futuras}

A partir del trabajo realizado, se identifican las siguientes líneas de desarrollo futuro:

\textbf{Implementación de Progressive Web App (PWA).} La conversión de la aplicación en una PWA, mediante la adición de un \textit{Service Worker} y un \textit{manifest} web, permitiría la instalación en dispositivos móviles y habilitaría un modo de funcionamiento \textit{offline} básico.

\textbf{Sistema de recomendaciones personalizado.} El historial de visionado acumulado por el usuario constituye un conjunto de datos con potencial para la implementación de algoritmos de recomendación colaborativa o basada en contenido, posiblemente mediante TensorFlow.js en el cliente o una función \textit{serverless}.

\textbf{Expansión de las funcionalidades sociales.} La comparación del historial entre usuarios, las recomendaciones de amigos y la visualización del perfil público de otros usuarios de Trakt son funcionalidades que la API de Trakt soporta pero no han sido implementadas en la versión actual.

\textbf{Backend propio y base de datos local.} A largo plazo, la eliminación de la dependencia total de Trakt.tv como fuente de verdad requeriría la implementación de un \textit{backend} propio con base de datos (PostgreSQL o MySQL) que almacenara el historial y las preferencias localmente.
